% archon:covers AlgebraicJacobian/RiemannRoch/RRFormula.lean

\chapter{The Riemann--Roch formula in genus zero (RR.2)}
\label{chap:RiemannRoch_RRFormula}

% SOURCE: Hartshorne, \emph{Algebraic Geometry}, IV \S1, pp.~294--297
% (read from references/hartshorne-algebraic-geometry.pdf, PDF pages
% 311--314). The verbatim quotes pasted in the declaration blocks below
% were copied character-by-character from the rendered page images of
% the scanned PDF (no text layer is present).

\section{Setup and motivation}

This chapter is the second of four sub-build chapters \texttt{RR.1}--\texttt{RR.4}
that close the project's headline RR bridge
\[
  \mathtt{genusZero\_curve\_iso\_P1}
  \,\colon\,
  \text{a smooth proper geometrically irreducible curve \(C / \bar k\) with
  \(g(C) = 0\) is isomorphic to \(\mathbb P^1_{\bar k}\).}
\]
The bridge consumed by
\cref{prop:rigidity_genus0_curve_to_AV} of
\texttt{AbelianVarietyRigidity.tex} is the Hartshorne~IV.1.3.5 chain:
the Weil-divisor group \(\mathrm{Div}(C)\) and the degree map $\deg \colon \mathrm{Div}(C)
\to \mathbb Z$ (\texttt{RR.1}, \cref{def:divisor_degree}); the
Riemann--Roch dimension formula in genus~\(0\)
\[
  \ell(D) \;=\; \deg(D) + 1 \qquad \text{for } \deg(D) \geq 0
\]
(this chapter, \texttt{RR.2}); the global-sections identification
\(\dim_{\bar k} H^0(C, \mathcal O_C(P)) = 2\) (\texttt{RR.3}); the
``rational curve \(\Rightarrow \cong \mathbb P^1\)'' classification via
\texttt{Proj.fromOfGlobalSections} (\texttt{RR.4}).

% NOTE (Serre duality vs. direct \(\chi\) — iter-173 sub-phase choice).
% This chapter proves the RR formula via the \emph{direct \(\chi\)
% inductive proof} (Hartshorne IV.1.3 step ``Next, let \(D\) be any
% divisor\ldots'', p.~296), \emph{not} via Serre duality. The reason is
% strictly Mathlib-side: Mathlib (snapshot \texttt{b80f227}) ships no
% dualizing sheaf \(\omega_X\) on a scheme and no Serre duality theorem
% \(H^i(X, \mathcal F) \cong H^{n-i}(X, \omega_X \otimes \mathcal F^\vee)^*\);
% routing through that machinery would force a multi-iter sub-build
% strictly heavier than the genus-\(0\) critical path requires. The
% direct \(\chi\) proof depends only on additivity of the Euler
% characteristic on a short exact sequence of coherent sheaves (a
% standard preadditive-category fact), the genus base case
% \(\chi(\mathcal O_C) = 1 - g\), and the closed-point skyscraper
% identity \(\chi(k(P)) = 1\).

% NOTE (Genus-\(0\) carve-out vs. general \(g\) — iter-173 sub-phase
% choice). The chapter restricts the headline theorem
% \cref{thm:euler_char_eq_deg_plus_one_minus_genus} to the genus-\(0\) specialisation
% \(\ell(D) = \deg(D) + 1\) for \(\deg(D) \geq 0\). The general-\(g\)
% statement \(\ell(D) - \ell(K - D) = \deg(D) + 1 - g\) is downstream of
% \texttt{RR.2} (it requires the canonical divisor \(K\) and Serre
% duality / the \(K - D\) pairing) and is \emph{not} needed for the
% genus-\(0\) critical path: \texttt{RR.3} consumes only
% \cref{thm:riemannRoch_genus_zero}, and \texttt{RR.4} consumes
% \texttt{RR.3}. The intermediate \(\chi\)-identity
% \cref{thm:euler_char_eq_deg_plus_one_minus_genus} is stated and
% sketched for general \(g\) (the inductive argument is genus-uniform);
% its full closure on a general-\(g\) curve waits for the project's
% subsequent Riemann--Roch chapter, but the proof body in
% \texttt{AlgebraicJacobian/RiemannRoch/RRFormula.lean} pins the
% \emph{statement} for general \(g\) and may close the body initially
% under a \(g = 0\) hypothesis.

\paragraph{Standing hypotheses.} Throughout this chapter, \(\bar k\) is an
algebraically closed field and \(C\) is a smooth proper geometrically
irreducible curve over \(\bar k\). Concretely, the Lean signatures
(\texttt{AlgebraicJacobian/RiemannRoch/RRFormula.lean}) carry the same
typeclass package as the rest of the project's curve API:
\begin{quote}
  \texttt{\{kbar : Type u\} [Field kbar] [IsAlgClosed kbar]\\
   (C : Over (Spec (.of kbar))) [IsProper C.hom]\\
   \phantom{}\quad[SmoothOfRelativeDimension 1 C.hom]\\
   \phantom{}\quad[GeometricallyIrreducible C.hom] [IsIntegral C.left]}
\end{quote}
The genus \(g(C)\) is the one defined in \cref{def:genus} of
\texttt{Genus.tex}, namely \(g(C) = \dim_{\bar k} H^1(C, \mathcal O_C)\)
via the project's \(\Module\bar k\)-flavoured sheaf cohomology
(\texttt{AlgebraicGeometry.genus}).

\section{The Euler characteristic of a coherent sheaf on a curve}

\begin{definition}


\leanok
[Euler characteristic of a coherent sheaf on a curve]
  \label{def:eulerChar_curve}
  \lean{AlgebraicGeometry.Scheme.eulerCharacteristic}
  \uses{def:Scheme_HModule}

  % SOURCE: [Hartshorne], IV.1, p.~295 (definition of \(\chi(\mathcal F)\)
  % inside the proof of Theorem~1.3) (read from
  % references/hartshorne-algebraic-geometry.pdf, PDF page 312).
  % SOURCE QUOTE: "Thus we have to show that for any \(D\),
  % \(\chi(\mathscr L(D)) = \deg D + 1 - g\),
  % where for any coherent sheaf \(\mathscr F\) on \(X\), \(\chi(\mathscr F)\)
  % is the Euler characteristic
  % \(\chi(\mathscr F) = \dim H^0(X, \mathscr F) - \dim H^1(X, \mathscr F)\)."
  \textit{Source: Hartshorne, IV.1, p.~295, the formula displayed
  inside the proof of Theorem~1.3.}
  Let \(C\) be a smooth proper curve over \(\bar k\) and let \(\mathcal F\)
  be a coherent sheaf on \(C\). The \emph{Euler characteristic} of
  \(\mathcal F\) is the integer
  \[
    \chi(\mathcal F) \;:=\;
    \sum_{i \geq 0} (-1)^i \,\dim_{\bar k} H^i(C, \mathcal F)
    \;=\;
    \dim_{\bar k} H^0(C, \mathcal F)
    \,-\,
    \dim_{\bar k} H^1(C, \mathcal F).
  \]
  The second equality is a curve-specific simplification: because
  \(\dim C = 1\), Grothendieck vanishing (Hartshorne~III.2.7) gives
  \(H^i(C, \mathcal F) = 0\) for \(i \geq 2\), so the alternating sum has
  only two terms. Coherence of \(\mathcal F\) on the proper $\bar
  k$-scheme $C$ guarantees that both $H^0$ and $H^1$ are
  finite-dimensional \(\bar k\)-vector spaces (Serre's coherent-cohomology
  finiteness theorem, the same finiteness backing
  \cref{def:l_invariant}), so the two terms are honest natural numbers and
  the difference is a well-defined integer.

  \paragraph{Lean signature scope.} The Lean definition
  \texttt{AlgebraicGeometry.Scheme.eulerCharacteristic} takes a coherent
  sheaf \(\mathcal F\) on \(C\) and returns an integer; the natural input
  shape is the project's \(\Module \bar k\)-valued sheaf-cohomology
  wrapper \(\HModule \bar k\) (\cref{def:Scheme_HModule}) applied to a
  \(\Module \bar k\)-flavoured carrier built from \(\mathcal F\), with
  \(\dim_{\bar k}\) extracted via \(\Module.\mathtt{finrank}\). The
  declaration is \texttt{noncomputable} (inherited from
  \(\Module.\mathtt{finrank}\) and from the \(\Ext\)-based cohomology
  construction).
\end{definition}

\section{The \(\ell\)-invariant of a Weil divisor}

\begin{definition}


\leanok
[\(\ell\)-invariant of a Weil divisor]
  \label{def:l_invariant}
  \lean{AlgebraicGeometry.Scheme.WeilDivisor.l}
  \uses{def:Scheme_HModule, def:codim1_cycles, def:sheafOf}

  % SOURCE: [Hartshorne], IV.1, p.~295 (definition of \(l(D)\) and of the
  % "dimension of \(|D|\)") (read from
  % references/hartshorne-algebraic-geometry.pdf, PDF page 312).
  % SOURCE QUOTE: "We denote \(\dim_k H^0(X, \mathscr L(D))\) by \(l(D)\),
  % so that the \emph{dimension of \(|D|\)} is \(l(D) - 1\). The number
  % \(l(D)\) is finite by (II, 5.19) or (III, 5.2)."
  \textit{Source: Hartshorne, IV.1, p.~295 (definition of \(l(D)\)).}
  Let \(C\) be a smooth proper curve over \(\bar k\) and let
  \(D \in \mathrm{Div}(C)\) be a Weil divisor (\cref{def:codim1_cycles}). The
  \emph{\(\ell\)-invariant of \(D\)} is the integer
  \[
    \ell(D) \;:=\; \dim_{\bar k} H^0(C, \mathcal O_C(D))
    \;\in\; \mathbb N,
  \]
  the \(\bar k\)-dimension of the space of global sections of the
  invertible sheaf \(\mathcal O_C(D)\) associated to \(D\). Finiteness of
  \(\ell(D)\) is a consequence of the coherent-cohomology finiteness
  theorem on a proper \(\bar k\)-scheme (Hartshorne~II.5.19 or
  III.5.2). The classical \emph{complete linear system \(|D|\)} has
  \(\bar k\)-dimension \(\ell(D) - 1\) as a projective space, but the
  chapter never uses the projective-space interpretation; \(\ell(D)\) is
  the only quantity it needs.

  \paragraph{Lean signature scope.} The Lean definition
  \texttt{AlgebraicGeometry.Scheme.WeilDivisor.l} accepts a divisor
  \(D : C.\mathtt{WeilDivisor}\) (\cref{def:codim1_cycles}) and returns
  \(\Module.\mathtt{finrank}\;\bar k\;(H^0(C, \mathcal O_C(D)))\),
  computed through the same \(\Module \bar k\)-valued sheaf-cohomology
  pipeline as \cref{def:eulerChar_curve}. The associated invertible
  sheaf \(\mathcal O_C(D)\) is the line bundle of \texttt{RR.3} (sibling
  chapter, to be added); the present chapter's prose and Lean
  signature reference \(\mathcal O_C(D)\) by name and consume only its
  \(H^0\) functor. The declaration is \texttt{noncomputable}.
\end{definition}

\section{The \(\chi\) identity: \(\chi(\mathcal O_C(D)) = \deg(D) + 1 - g\)}

\subsection*{Substrate lemmas for the proof of the \(\chi\)-identity}

The proof of \cref{thm:euler_char_eq_deg_plus_one_minus_genus} below
decomposes into three substantive substrate lemmas, factored out of the
Hartshorne~IV.1.3 inductive argument. The first
(\cref{lem:euler_char_shortExact_add}, \emph{additivity}) is the
classical statement that \(\chi\) descends to the Grothendieck group
of \(\mathbf{Coh}(C)\); the second
(\cref{lem:euler_char_iso}, \emph{iso-invariance}) is the
elementary fact that \(\chi\) is well defined on isomorphism classes
of sheaves; the third
(\cref{lem:euler_char_skyscraperSheaf}, \emph{skyscraper}) computes \(\chi\)
on a closed-point skyscraper sheaf \(k(P)\). Pinning each as its own
Lean declaration lets the prover attack the three substrate gaps
independently of the parent inductive argument.

\begin{lemma}\leanok
[Additivity of the Euler characteristic on a short exact sequence]
  \label{lem:euler_char_shortExact_add}
  \lean{AlgebraicGeometry.Scheme.eulerCharacteristic_shortExact_add}
  \uses{def:eulerChar_curve, def:Scheme_HModule}
  % SOURCE: [Hartshorne], IV.1, p.~296 (inside the proof of Theorem~1.3,
  % the citation ``III, Ex.~5.1'' for the additivity of \(\chi\) on a
  % short exact sequence) (read from
  % references/hartshorne-algebraic-geometry.pdf, PDF page 313).
  % SOURCE QUOTE: "Now the Euler characteristic is additive on short
  % exact sequences (III, Ex.~5.1), and \(\chi(k(P)) = 1\), so we have
  % \(\chi(\mathscr L(D + P)) = \chi(\mathscr L(D)) + 1\)."
  \textit{Source: Hartshorne, IV.1 p.~296 (citation of III, Ex.~5.1
  inside the proof of Theorem~1.3).}
  Let \(C\) be a smooth proper geometrically irreducible curve over
  \(\bar k\), and let
  \[
    0 \to \mathcal F_1 \to \mathcal F_2 \to \mathcal F_3 \to 0
  \]
  be a short exact sequence in the category of \(\Module \bar k\)-valued
  sheaves on \(C\). Then
  \[
    \chi(\mathcal F_2) \;=\; \chi(\mathcal F_1) + \chi(\mathcal F_3).
  \]
\end{lemma}

\begin{proof}
  \uses{def:eulerChar_curve}
  Apply the project-side \(\Module \bar k\)-flavoured covariant Ext-LES
  (Mathlib's \texttt{Abelian.Ext.covariantSequence}) to the constant-sheaf
  test object \(L = (\mathrm{constantSheaf}\;J\;\_)(\Module.\mathtt{of}\;
  \bar k\;\bar k)\) and the short exact sequence
  \(0 \to \mathcal F_1 \to \mathcal F_2 \to \mathcal F_3 \to 0\). The
  resulting nine-term sequence
  \[
    H^0(\mathcal F_1) \to H^0(\mathcal F_2) \to H^0(\mathcal F_3)
    \to H^1(\mathcal F_1) \to H^1(\mathcal F_2) \to H^1(\mathcal F_3)
    \to H^2(\mathcal F_1) \to H^2(\mathcal F_2) \to H^2(\mathcal F_3)
  \]
  is exact in the category of \(\bar k\)-vector spaces. Grothendieck
  vanishing \(H^i(C, \mathcal F) = 0\) for \(i \geq 2\) on the
  one-dimensional scheme \(C\) (Hartshorne~III.2.7) truncates the
  sequence to a six-term exact sequence
  \[
    0 \to K_0 \to H^0(\mathcal F_1) \to H^0(\mathcal F_2) \to
    H^0(\mathcal F_3) \to H^1(\mathcal F_1) \to H^1(\mathcal F_2)
    \to H^1(\mathcal F_3) \to 0
  \]
  (with \(K_0 = 0\) the kernel of the first map, kept here only as a
  bookkeeping placeholder). Rank-counting on a six-term exact sequence
  of finite-dimensional \(\bar k\)-vector spaces — iterated applications
  of \(\mathtt{Submodule.finrank\_quotient\_add\_finrank}\) at each
  intermediate kernel/image quotient — yields the alternating identity
  \[
    \dim H^0(\mathcal F_1) - \dim H^0(\mathcal F_2) + \dim H^0(\mathcal F_3)
    - \dim H^1(\mathcal F_1) + \dim H^1(\mathcal F_2) - \dim H^1(\mathcal F_3)
    \;=\; 0,
  \]
  which rearranges, using the definition
  \(\chi(\mathcal F) = \dim H^0 - \dim H^1\)
  (\cref{def:eulerChar_curve}), to
  \(\chi(\mathcal F_2) = \chi(\mathcal F_1) + \chi(\mathcal F_3)\).

  % NOTE: gated on project-side LES carrier + finiteness machinery.
  % Mathlib `b80f227` ships `Abelian.Ext.covariantSequence` but not (a)
  % the `ModuleCat kbar`-flavoured LES carrier specialised to the
  % project's sheaf-cohomology pipeline `Scheme.HModule`, (b) Grothendieck
  % vanishing `H^i(C, F) = 0` for `i ≥ 2` at the `Ext`-cohomology level
  % consumed by `HModule`, nor (c) the six-term rank-counting bridge
  % composed of `Submodule.finrank_quotient_add_finrank` invocations.
  % The Lean helper `eulerCharacteristic_shortExact_add` therefore
  % carries a Tier-3 honest typed sorry; its closure lands once the
  % three substrate pieces are in place. Off the genus-0 critical path
  % (the parent theorem `thm:euler_char_eq_deg_plus_one_minus_genus`
  % is the only consumer in `RR.2`).
\end{proof}

\begin{lemma}\leanok
[Iso-invariance of the Euler characteristic]
  \label{lem:euler_char_iso}
  \lean{AlgebraicGeometry.Scheme.eulerCharacteristic_iso}
  \uses{def:eulerChar_curve}
  \textit{Project-bespoke decomposition} (Archon-original packaging of
  the Mathlib \(\Ext\)-postcomposition API; no external source named).
  Let \(C\) be a smooth proper geometrically irreducible curve over
  \(\bar k\) and let \(\mathcal F, \mathcal G\) be \(\Module \bar k\)-valued
  sheaves on \(C\) with an isomorphism \(\mathcal F \cong \mathcal G\).
  Then
  \[
    \chi(\mathcal F) \;=\; \chi(\mathcal G).
  \]
\end{lemma}

\begin{proof}
  \uses{def:eulerChar_curve}
  \leanok
  An isomorphism \(e \colon \mathcal F \xrightarrow{\;\cong\;} \mathcal G\)
  induces, in each cohomological degree \(n \in \mathbb N\), a
  \(\bar k\)-linear equivalence
  \[
    H^n(C, \mathcal F) \;\xrightarrow{\;\cong\;}\; H^n(C, \mathcal G)
  \]
  by post-composition with the Ext-degree-zero classes
  \(\mathtt{Ext.mk}_0(e.\mathtt{hom})\) and
  \(\mathtt{Ext.mk}_0(e.\mathtt{inv})\). The mutual-inverse identities
  for this pair of post-compositions reduce — via the associativity
  lemma \(\mathtt{Ext.comp\_assoc\_of\_third\_deg\_zero}\), the
  composition lemma \(\mathtt{Ext.mk}_0\mathtt{\_comp\_mk}_0\), and the
  iso identities \(e.\mathtt{hom} \circ e.\mathtt{inv} = \mathrm{id}\),
  \(e.\mathtt{inv} \circ e.\mathtt{hom} = \mathrm{id}\) — to the identity
  map on \(\Ext\) elements. Specialising to \(n = 0\) and \(n = 1\) and
  applying \(\mathtt{LinearEquiv.finrank\_eq}\) at each of those two
  degrees gives
  \(\dim H^n(C, \mathcal F) = \dim H^n(C, \mathcal G)\) for
  \(n \in \{0, 1\}\). Unfolding the definition
  \(\chi(\mathcal F) = \dim H^0 - \dim H^1\)
  (\cref{def:eulerChar_curve}) closes the goal.
\end{proof}

\begin{lemma}\leanok
[H\(^0\) of a closed-point skyscraper sheaf has dimension 1]
  \label{lem:H0_skyscraperSheaf_finrank_eq_one}
  \lean{AlgebraicGeometry.Scheme.H0_skyscraperSheaf_finrank_eq_one}
  \uses{def:Scheme_HModule_zero_linearEquiv,
        def:Scheme_constantSheafGammaHom_linearEquiv,
        thm:Scheme_SheafGammaObj_linearEquiv_top}
  Iter-188 Lane H named substrate. Let \(C\) be a smooth proper
  geometrically irreducible curve over \(\bar k\) and \(P \in C\) a closed
  point. Then
  \[
    \dim_{\bar k}\, H^0(C, \mathtt{skyscraperSheaf}\, P\,
      (\Module.\mathtt{of}\, \bar k\, \bar k)) \;=\; 1.
  \]
  Closed axiom-clean iter-188 via the three-step \(\bar k\)-linear-
  equivalence chain \texttt{HModule\_zero\_linearEquiv}
  \(\to\) \texttt{constantSheafGammaHom\_linearEquiv}
  \(\to\) \texttt{SheafGammaObj\_linearEquiv\_top}, followed by
  \texttt{LinearEquiv.finrank\_eq} and \texttt{Module.finrank\_self}
  on the dite-resolved skyscraper presheaf evaluation at \(\top\).
\end{lemma}

\begin{lemma}

\leanok
[H\(^1\) of a closed-point skyscraper sheaf vanishes]
  \label{lem:H1_skyscraperSheaf_finrank_eq_zero}
  \lean{AlgebraicGeometry.Scheme.H1_skyscraperSheaf_finrank_eq_zero}
  \uses{lem:H1_skyscraperSheaf_finrank_eq_zero_main}
  Iter-188 Lane H named substrate (typed-sorry blocker; committed
  RR.2.H\(^1\) STRATEGY sub-phase). Let \(C\) be as above. Then
  \[
    \dim_{\bar k}\, H^1(C, \mathtt{skyscraperSheaf}\, P\,
      (\Module.\mathtt{of}\, \bar k\, \bar k)) \;=\; 0.
  \]
  Mathematically: the closed-point skyscraper sheaf is flasque (any
  pushforward from a one-point space is flasque; Hartshorne~III.2.5(c)),
  and flasque sheaves have vanishing higher cohomology on any
  topological space (Hartshorne~III.2.5).
  Gated on either (a) upstream Mathlib
  \texttt{Sheaf.flasque\_vanishing\_cohomology} at the
  \texttt{Ext}-cohomology level for sheaves valued in an abelian
  category with enough injectives, or (b) a project-side
  flasque-cohomology bridge (~150-300 LOC; the committed
  \cref{chap:RR_H1Vanishing}-style sub-build).
\end{lemma}

\begin{lemma}\leanok
[Euler characteristic of a closed-point skyscraper sheaf]
  \label{lem:euler_char_skyscraperSheaf}
  \lean{AlgebraicGeometry.Scheme.eulerCharacteristic_skyscraperSheaf}
  \uses{def:eulerChar_curve,
        lem:H0_skyscraperSheaf_finrank_eq_one,
        lem:H1_skyscraperSheaf_finrank_eq_zero}
  % SOURCE: [Hartshorne], IV.1, p.~296 (the \(\chi(k(P)) = 1\) input
  % cited inline inside the proof of Theorem~1.3) (read from
  % references/hartshorne-algebraic-geometry.pdf, PDF page 313).
  % SOURCE QUOTE: "Now the Euler characteristic is additive on short
  % exact sequences (III, Ex.~5.1), and \(\chi(k(P)) = 1\), so we have
  % \(\chi(\mathscr L(D + P)) = \chi(\mathscr L(D)) + 1\)."
  \textit{Source: Hartshorne, IV.1 p.~296 (the ``\(\chi(k(P)) = 1\)''
  input inside the proof of Theorem~1.3).}
  Let \(C\) be a smooth proper geometrically irreducible curve over
  \(\bar k\) and let \(P \in C\) be a closed point. Write \(k(P)\) for
  the closed-point skyscraper sheaf
  \[
    k(P) \;:=\; \mathtt{skyscraperSheaf}\;P\;
    (\Module.\mathtt{of}\;\bar k\;\bar k),
  \]
  the \(\Module \bar k\)-valued sheaf on \(C\) with stalk \(\bar k\) at
  \(P\) and zero stalks elsewhere. Then
  \[
    \chi(k(P)) \;=\; 1.
  \]
\end{lemma}

\begin{proof}
  \uses{def:eulerChar_curve,
        lem:H0_skyscraperSheaf_finrank_eq_one,
        lem:H1_skyscraperSheaf_finrank_eq_zero}
  \leanok
  By the definition of \(\chi\)
  (\cref{def:eulerChar_curve}),
  \[
    \chi(k(P)) \;=\; \dim_{\bar k} H^0(C, k(P))
    \;-\; \dim_{\bar k} H^1(C, k(P)).
  \]
  The closure decomposes cleanly into two independent half-statements
  (\cref{lem:H0_skyscraperSheaf_finrank_eq_one} and
  \cref{lem:H1_skyscraperSheaf_finrank_eq_zero}), attacked separately
  by the prover.

  \emph{H\(^0\) half} — \(\dim_{\bar k} H^0(C, k(P)) = 1\).
  Compose the project's three-step \(\bar k\)-linear-equivalence chain
  \[
    H^0(C, k(P))
    \;\xrightarrow{\;\cong\;}\;
    \mathrm{Hom}_{\mathrm{Sh}}\!\bigl(
      (\mathrm{constantSheaf}\;J\;\_)(\Module.\mathtt{of}\;\bar k\;\bar k),\,
      k(P)
    \bigr)
    \;\xrightarrow{\;\cong\;}\;
    \Gamma(C, k(P))
    \;\xrightarrow{\;\cong\;}\;
    (k(P))(C),
  \]
  realised by the project lemmas
  \texttt{Scheme.HModule\_zero\_linearEquiv},
  \texttt{Scheme.constantSheafGammaHom\_linearEquiv}, and
  \texttt{Scheme.SheafGammaObj\_linearEquiv\_top} (the last step
  identifies the abstract \(\Gamma\) functor with the underlying
  presheaf evaluation at the top open \(\top\), i.e.\ at the whole
  underlying space of \(C\)). Applying
  \(\mathtt{LinearEquiv.finrank\_eq}\) along this chain reduces the
  goal to
  \(\dim_{\bar k} (k(P))(C) = 1\). The closed point
  \(P\) lies in the top open of \(C\), so the skyscraper presheaf
  evaluates at the top open to its stalk
  \(\Module.\mathtt{of}\;\bar k\;\bar k\) by definition, and
  \(\dim_{\bar k} \bar k = 1\) closes the half. This half is
  axiom-clean closure projection (\(\sim 30\text{--}60\) LOC) against
  the existing project chain; it is the natural target for the
  iteration that follows the present blueprint expansion.

  \emph{H\(^1\) half} — \(\dim_{\bar k} H^1(C, k(P)) = 0\).
  Mathematically, the closed-point skyscraper sheaf is flasque (any
  pushforward of a sheaf from a one-point space is flasque;
  Hartshorne~III.2.5(c)), and flasque sheaves have vanishing higher
  cohomology in degrees \(\geq 1\) on any topological space
  (Hartshorne~III.2.5). Hence \(H^1(C, k(P)) = 0\) as a
  \(\bar k\)-vector space.

  % NOTE: gated on Mathlib flasque cohomology.
  % Mathlib `b80f227` does NOT ship an `H^n(X, F) = 0 for flasque F`
  % lemma at the `ModuleCat`-flavoured `Abelian.Ext`-cohomology level
  % used by the project's `Scheme.HModule kbar` pipeline. The closure
  % depends on either (a) an upstream Mathlib
  % `Sheaf.flasque_vanishing_cohomology` contribution at the
  % `Ext`-cohomology level for sheaves with values in an abelian
  % category with enough injectives, or (b) a project-side flasque-
  % cohomology bridge (~150-300 LOC, off the genus-0 critical path).
  % Until one of the two lands, the H¹ half is a typed-sorry blocker
  % independent of the H⁰ half — i.e. the prover can attack the H⁰
  % half against the existing project chain without waiting for the
  % H¹ machinery, and the present skyscraper lemma waits for the H¹
  % gap to clear.

  Combining the two halves through the definition of \(\chi\),
  \(\chi(k(P)) = 1 - 0 = 1\), closes the lemma.
\end{proof}

\begin{lemma}[Global sections of the structure sheaf are one-dimensional]
  \label{lem:finrank_H0_toModuleKSheaf_eq_one}
  \lean{AlgebraicGeometry.Scheme.finrank_H0_toModuleKSheaf_eq_one}
  \uses{def:Scheme_HModule, def:Scheme_toModuleKSheaf}
  \textit{Specialisation of Hartshorne~I.3.4 (``\(H^0(X, \mathcal O_X) = k\)
  for any projective variety over an algebraically closed field'') to the
  project's \(\Module \bar k\)-flavoured cohomology pipeline.}
  Let \(C\) be a smooth proper geometrically irreducible curve over
  \(\bar k\). The zeroth cohomology of the structure-sheaf carrier
  \(\toModuleKSheaf C\) is one-dimensional over \(\bar k\):
  \[
    \dim_{\bar k} H^0(C, \mathcal O_C)
    \;=\; \dim_{\bar k}\, \HModule\,\bar k\,(\toModuleKSheaf C)\,0
    \;=\; 1.
  \]
  The proof identifies \(H^0(C, \mathcal O_C)\) with the global-sections
  ring \(\Gamma(C, \mathcal O_C) = \mathcal O_C(\top)\) through the
  project's \(\bar k\)-linear-equivalence chain
  (\cref{def:Scheme_HModule}). On the proper \(\bar k\)-scheme \(C\) this
  ring is a finite-dimensional \(\bar k\)-algebra; integrality of \(C\)
  makes it a domain; and a finite \(\bar k\)-algebra that is an integral
  domain over the algebraically closed field \(\bar k\) is \(\bar k\)
  itself, so its \(\bar k\)-dimension is \(1\).
\end{lemma}

\begin{proof}
  Proved directly in Lean.
\end{proof}

\begin{lemma}[\(\chi\)-additivity across a short exact sequence with skyscraper cokernel]
  \label{lem:euler_char_of_shortExact_skyscraper}
  \lean{AlgebraicGeometry.Scheme.eulerCharacteristic_of_shortExact_skyscraper}
  \uses{def:eulerChar_curve, lem:euler_char_shortExact_add,
        lem:euler_char_iso, lem:euler_char_skyscraperSheaf}
  \textit{Project-bespoke packaging of the Hartshorne~IV.1.3 inductive
  step, combining the three substrate lemmas above.}
  Let \(C\) be a smooth proper geometrically irreducible curve over
  \(\bar k\), let
  \[
    0 \to \mathcal F_1 \to \mathcal F_2 \to \mathcal F_3 \to 0
  \]
  be a short exact sequence of \(\Module \bar k\)-valued sheaves on
  \(C\), and let \(P \in C\) be a closed point such that the cokernel
  term \(\mathcal F_3\) is isomorphic to the skyscraper sheaf
  \(k(P) = \mathtt{skyscraperSheaf}\;P\;(\Module.\mathtt{of}\;\bar k\;\bar k)\).
  Then
  \[
    \chi(\mathcal F_2) \;=\; \chi(\mathcal F_1) + 1.
  \]
\end{lemma}

\begin{proof}
  \uses{lem:euler_char_shortExact_add, lem:euler_char_iso,
        lem:euler_char_skyscraperSheaf}
  Additivity of the Euler characteristic across the short exact sequence
  (\cref{lem:euler_char_shortExact_add}) gives
  \(\chi(\mathcal F_2) = \chi(\mathcal F_1) + \chi(\mathcal F_3)\).
  Iso-invariance of \(\chi\) (\cref{lem:euler_char_iso}) applied to the
  isomorphism \(\mathcal F_3 \cong k(P)\) rewrites the last term as
  \(\chi(k(P))\), and the skyscraper computation
  (\cref{lem:euler_char_skyscraperSheaf}) evaluates \(\chi(k(P)) = 1\).
  Combining the three identities yields
  \(\chi(\mathcal F_2) = \chi(\mathcal F_1) + 1\).
\end{proof}

\subsection*{Inductive scaffolding via the divisor line bundle}

The base case and the single-point inductive step of the
\(\chi\)-identity \cref{thm:euler_char_eq_deg_plus_one_minus_genus} are
pinned as three named Lean lemmas. Each references the line bundle
\(\mathcal O_C(D)\) of \texttt{RR.3} (\cref{def:sheafOf}) and so carries a
typed-sorry transitivity inherited from that sibling chapter; the
arithmetic scaffolding around \(\chi\) is, however, honest.

\begin{lemma}[Base case of the \(\chi\)-identity]
  \label{lem:euler_char_sheafOf_zero}
  \lean{AlgebraicGeometry.Scheme.eulerCharacteristic_sheafOf_zero}
  \uses{def:eulerChar_curve, lem:finrank_H0_toModuleKSheaf_eq_one,
        lem:sheafOf_zero, def:genus}
  \textit{Base case of Hartshorne~IV.1.3 (the \(D = 0\) step).}
  Let \(C\) be a smooth proper geometrically irreducible curve over
  \(\bar k\). The Euler characteristic of the line bundle
  \(\mathcal O_C(0)\) of the zero divisor is
  \[
    \chi(\mathcal O_C(0)) \;=\; 1 - g(C).
  \]
\end{lemma}

\begin{proof}
  \uses{lem:finrank_H0_toModuleKSheaf_eq_one, lem:sheafOf_zero, def:genus}
  The line bundle of the zero divisor is the structure-sheaf carrier,
  \(\mathcal O_C(0) = \toModuleKSheaf C\) (\cref{lem:sheafOf_zero}).
  Unfolding the two-term Euler characteristic
  \(\chi = \dim_{\bar k} H^0 - \dim_{\bar k} H^1\)
  (\cref{def:eulerChar_curve}) reduces the claim to
  \[
    \dim_{\bar k} H^0(C, \mathcal O_C)
    - \dim_{\bar k} H^1(C, \mathcal O_C) \;=\; 1 - g(C).
  \]
  The first term equals \(1\)
  (\cref{lem:finrank_H0_toModuleKSheaf_eq_one}), and the second term is by
  definition the genus, \(\dim_{\bar k} H^1(C, \mathcal O_C) = g(C)\)
  (\cref{def:genus}). Hence \(\chi(\mathcal O_C(0)) = 1 - g(C)\).
\end{proof}

\begin{lemma}[Unit inductive step of the \(\chi\)-identity]
  \label{lem:euler_char_sheafOf_succ}
  \lean{AlgebraicGeometry.Scheme.eulerCharacteristic_sheafOf_succ}
  \uses{def:eulerChar_curve, lem:euler_char_of_shortExact_skyscraper,
        lem:sheafOf_ses_single_add, def:sheafOf}
  \textit{Single-point inductive step of Hartshorne~IV.1.3 (the
  ``\(D \leftrightsquigarrow D + P\)'' step, additive direction).}
  Let \(C\) be as above, let \(D \in \mathrm{Div}(C)\), and let \(Y \in C\)
  be a prime (closed-point) divisor. Adding the single point \(Y\) raises
  the Euler characteristic by one:
  \[
    \chi(\mathcal O_C([Y] + D)) \;=\; \chi(\mathcal O_C(D)) + 1.
  \]
\end{lemma}

\begin{proof}
  \uses{lem:euler_char_of_shortExact_skyscraper, lem:sheafOf_ses_single_add}
  The closed-point short exact sequence of \texttt{RR.3}
  (\cref{lem:sheafOf_ses_single_add}) provides a short exact sequence
  \[
    0 \to \mathcal O_C(D) \to \mathcal O_C([Y] + D) \to k(Y) \to 0
  \]
  of \(\Module \bar k\)-valued sheaves whose cokernel is the closed-point
  skyscraper \(k(Y)\). Applying the packaged \(\chi\)-additivity with
  skyscraper cokernel (\cref{lem:euler_char_of_shortExact_skyscraper}) to
  this sequence yields
  \(\chi(\mathcal O_C([Y] + D)) = \chi(\mathcal O_C(D)) + 1\).
\end{proof}

\begin{lemma}[General single-point inductive step of the \(\chi\)-identity]
  \label{lem:euler_char_sheafOf_single_add}
  \lean{AlgebraicGeometry.Scheme.eulerCharacteristic_sheafOf_single_add}
  \uses{def:eulerChar_curve, lem:euler_char_sheafOf_succ, def:sheafOf}
  \textit{Iterated single-point step of Hartshorne~IV.1.3, both signs of
  the multiplicity.}
  Let \(C\) be as above, let \(D \in \mathrm{Div}(C)\), let \(Y \in C\) be
  a prime divisor, and let \(n \in \mathbb Z\). Then
  \[
    \chi(\mathcal O_C(n[Y] + D)) \;=\; \chi(\mathcal O_C(D)) + n.
  \]
\end{lemma}

\begin{proof}
  \uses{lem:euler_char_sheafOf_succ}
  Induction on \(n \in \mathbb Z\). For \(n = 0\) the divisor \(n[Y]\)
  vanishes and the identity is trivial. For the positive direction,
  \((n+1)[Y] + D = [Y] + (n[Y] + D)\), so the unit step
  (\cref{lem:euler_char_sheafOf_succ}) and the inductive hypothesis give
  \[
    \chi(\mathcal O_C((n+1)[Y] + D))
    \;=\; \chi(\mathcal O_C(n[Y] + D)) + 1
    \;=\; \chi(\mathcal O_C(D)) + (n+1).
  \]
  The negative direction runs the same unit step backwards: applying
  \cref{lem:euler_char_sheafOf_succ} to the divisor \((n-1)[Y] + D\) and
  rearranging gives
  \(\chi(\mathcal O_C((n-1)[Y] + D)) = \chi(\mathcal O_C(D)) + (n-1)\),
  which is the inductive hypothesis transported one step in the negative
  direction.
\end{proof}

\subsection*{The main \(\chi\)-identity}

\begin{theorem}


\leanok
[Euler-characteristic identity for \(\mathcal O_C(D)\)]
  \label{thm:euler_char_eq_deg_plus_one_minus_genus}
  \lean{AlgebraicGeometry.Scheme.eulerCharacteristic_eq_degree_plus_one_minus_genus}
  \uses{def:eulerChar_curve, def:l_invariant,
        lem:euler_char_shortExact_add, lem:euler_char_iso,
        lem:euler_char_skyscraperSheaf, def:genus,
        def:divisor_degree, def:divisor_closed_point, def:codim1_cycles}

  % SOURCE: [Hartshorne], IV.1, Theorem~1.3 (the reduction inside the
  % proof: ``Thus we have to show that for any \(D\), $\chi(\mathscr L(D))
  % = \deg D + 1 - g$'') (read from
  % references/hartshorne-algebraic-geometry.pdf, PDF page 312).
  % SOURCE QUOTE: "Theorem 1.3 (Riemann--Roch). \emph{Let \(D\) be a
  % divisor on a curve \(X\) of genus \(g\). Then}
  % \(l(D) - l(K - D) = \deg D + 1 - g\)."
  % SOURCE QUOTE: "Thus we have to show that for any \(D\),
  % \(\chi(\mathscr L(D)) = \deg D + 1 - g\),
  % where for any coherent sheaf \(\mathscr F\) on \(X\), \(\chi(\mathscr F)\)
  % is the Euler characteristic
  % \(\chi(\mathscr F) = \dim H^0(X, \mathscr F) - \dim H^1(X, \mathscr F)\)."
  \textit{Source: Hartshorne, IV.1, Theorem~1.3, p.~295.}
  Let \(C\) be a smooth proper geometrically irreducible curve over
  \(\bar k\) of genus \(g = g(C)\), and let \(D \in \mathrm{Div}(C)\). Then
  \[
    \chi(\mathcal O_C(D)) \;=\; \deg(D) + 1 - g.
  \]
\end{theorem}

% SOURCE QUOTE PROOF: "First we consider the case \(D = 0\). Then our
% formula says
% \(\dim H^0(X, \mathscr O_X) - \dim H^1(X, \mathscr O_X) = 1 - g\),
% because \(H^0(X, \mathscr O_X) = k\) for any projective variety
% (I, 3.4), and \(\dim H^1(X, \mathscr O_X) = g\) by (1.1).
% Next, let \(D\) be any divisor, and let \(P\) be any point. We will show
% that the formula is true for \(D\) if and only if it is true for
% \(D + P\). Since any divisor can be reached from \(0\) in a finite number
% of steps by adding or subtracting a point each time, this will show
% the result holds for all \(D\).
% We consider \(P\) as a closed subscheme of \(X\). Its structure sheaf is
% the skyscraper sheaf \(k\) sitting at the point \(P\), which we denote by
% \(k(P)\), and its ideal sheaf is \(\mathscr L(-P)\) by (II, 6.18).
% Therefore we have an exact sequence
% \(0 \to \mathscr L(-P) \to \mathscr O_X \to k(P) \to 0\).
% Tensoring with \(\mathscr L(D + P)\) we get
% \(0 \to \mathscr L(D) \to \mathscr L(D + P) \to k(P) \to 0\),
% (since \(\mathscr L(P)\) is locally free of rank 1, tensoring by it
% does not affect the sheaf \(k(P)\)). Now the Euler characteristic is
% additive on short exact sequences (III, Ex. 5.1), and $\chi(k(P))
% = 1$, so we have
% \(\chi(\mathscr L(D + P)) = \chi(\mathscr L(D)) + 1\).
% On the other hand, \(\deg D + P = \deg D + 1\), so our formula is true
% for \(D\) if and only if it is true for \(D + P\), as required."
\begin{proof}
  \leanok
  \uses{def:eulerChar_curve, def:l_invariant,
    lem:euler_char_shortExact_add, lem:euler_char_iso,
    lem:euler_char_skyscraperSheaf}

  The proof is by induction on the structure of \(\mathrm{Div}(C)\) as the free
  abelian group on closed points (\cref{def:codim1_cycles},
  \cref{def:divisor_closed_point}). Both sides of the asserted
  equation are additive in \(D\) (the right-hand side by linearity of
  \(\deg\), \cref{thm:divisor_degree_hom}; the left-hand side by the
  inductive step proved below), so it suffices to verify the formula
  for \(D = 0\) and to show that it transports across the elementary
  modification \(D \mapsto D + [P]\) at any closed point \(P \in C\).

  \emph{Base case \(D = 0\).} The line bundle \(\mathcal O_C(0)\) is the
  structure sheaf \(\mathcal O_C\), and the formula
  \(\chi(\mathcal O_C) = 1 - g\) reads
  \[
    \dim_{\bar k} H^0(C, \mathcal O_C)
    - \dim_{\bar k} H^1(C, \mathcal O_C)
    \;=\; 1 - g.
  \]
  The second term is the definition of the genus,
  \(\dim_{\bar k} H^1(C, \mathcal O_C) = g\) (\cref{def:genus}). The
  first term is \(\dim_{\bar k} H^0(C, \mathcal O_C) = 1\): \(C\) is a
  proper geometrically irreducible curve over the algebraically
  closed field \(\bar k\), so its connected components agree with its
  geometrically connected components, and $H^0(C, \mathcal O_C)
  \cong \bar k$ as $\bar k$-vector spaces (the Hartshorne~I.3.4
  statement that ``\(H^0(X, \mathcal O_X) = k\) for any projective
  variety'' over an algebraically closed field; the Mathlib home is
  the project's \(H^0\)-bridge in
  \texttt{Cohomology\_StructureSheafModuleK.tex}). Substituting these
  two values gives \(1 - g = 1 - g\), as required.

  \emph{Inductive step (Hartshorne's ``\(D \leftrightsquigarrow D + P\)''
  step).} Let \(P \in C\) be a closed point. We follow Hartshorne's
  argument: view \(P\) as a closed subscheme of \(C\); its structure
  sheaf is the skyscraper sheaf \(k(P) \cong \bar k\) supported at \(P\),
  and its sheaf of ideals is \(\mathcal O_C(-[P]) \subset \mathcal O_C\)
  (the locally principal ideal sheaf cutting out \(P\) in \(C\)). Hence
  the short exact sequence
  \[
    0 \to \mathcal O_C(-[P]) \to \mathcal O_C \to k(P) \to 0
  \]
  in \(\mathbf{Coh}(C)\) (Hartshorne~II.6.18). Tensoring with the
  invertible sheaf \(\mathcal O_C(D + [P])\) (a locally free
  \(\mathcal O_C\)-module of rank \(1\), so tensoring preserves exactness
  and leaves the skyscraper sheaf \(k(P)\) unchanged) yields
  \[
    0 \to \mathcal O_C(D)
    \to \mathcal O_C(D + [P])
    \to k(P) \to 0.
  \]
  (Strictly: tensoring \(\mathcal O_C(-[P]) \subset \mathcal O_C\) with
  \(\mathcal O_C(D + [P])\) produces a sheaf isomorphic — but not
  syntactically equal — to \(\mathcal O_C(D)\), and the skyscraper
  factor \(k(P)\) is preserved up to isomorphism because
  \(\mathcal O_C(D + [P])\) is locally free of rank \(1\). We transport
  along this isomorphism via iso-invariance of \(\chi\),
  \cref{lem:euler_char_iso}.)
  Additivity of \(\chi\) on the short exact sequence
  (\cref{lem:euler_char_shortExact_add}, the project-side
  factoring of the ``Now the Euler characteristic is additive on short
  exact sequences (III, Ex.~5.1)'' input of Hartshorne~IV.1.3) gives
  \[
    \chi(\mathcal O_C(D + [P]))
    \;=\;
    \chi(\mathcal O_C(D)) + \chi(k(P)).
  \]
  The skyscraper-sheaf input
  \(\chi(k(P)) = 1\) is the standalone substrate lemma
  \cref{lem:euler_char_skyscraperSheaf} (factored from the
  same Hartshorne~IV.1.3 line); its proof splits into the
  H\(^0\) identification \(\dim_{\bar k} H^0(C, k(P)) = 1\) and the
  H\(^1\) vanishing \(\dim_{\bar k} H^1(C, k(P)) = 0\), pursued
  independently by the prover (see the lemma's two-paragraph proof
  body). Therefore
  \[
    \chi(\mathcal O_C(D + [P])) \;=\; \chi(\mathcal O_C(D)) + 1.
  \]
  On the other hand, \(\deg([P]) = 1\) (every closed point contributes
  \(1\) to the bare-sum degree of \cref{def:divisor_degree}), so by
  \cref{thm:divisor_degree_hom}
  \[
    \deg(D + [P]) + 1 - g \;=\; \deg(D) + 1 - g + 1.
  \]
  Comparing the two displayed identities, the formula
  \(\chi(\mathcal O_C(D)) = \deg(D) + 1 - g\) holds for \(D\) if and only
  if it holds for \(D + [P]\). Since the closed-point divisors \([P]\)
  generate \(\mathrm{Div}(C)\) as a free abelian group, the formula propagates
  from \(D = 0\) to every \(D \in \mathrm{Div}(C)\) by induction on the (signed)
  number of generators in the expression of \(D\).

  \emph{Lean reference note on the three substrate lemmas.} The
  inductive step of this proof consumes three substantive substrate
  lemmas, each pinned to its own Lean declaration in
  \texttt{AlgebraicJacobian/RiemannRoch/RRFormula.lean}: the additivity
  \cref{lem:euler_char_shortExact_add} (\(\chi\) on a short
  exact sequence, gated on the project-side
  \(\Module \bar k\)-flavoured Ext-LES carrier + Grothendieck
  vanishing + six-term rank-counting), the iso-invariance
  \cref{lem:euler_char_iso} (\(\chi\) on
  isomorphism classes, axiom-clean against Mathlib's
  \(\Ext\)-postcomposition API; closed iter-186), and the skyscraper
  identity \cref{lem:euler_char_skyscraperSheaf}
  (\(\chi(k(P)) = 1\), splitting into an H\(^0\) half closed against
  the project's existing \(H^0\)-bridge and an H\(^1\) half gated on
  Mathlib flasque-cohomology — see the lemma's annotated proof body).
  Mathlib \texttt{b80f227} provides the Euler-characteristic API
  \texttt{Mathlib.Algebra.Homology.EulerCharacteristic}
  (\texttt{HomologicalComplex.eulerChar},
  \texttt{GradedObject.eulerChar}) but does \emph{not} ship the
  specific additivity statement
  \texttt{CategoryTheory.ShortExact.eulerChar\_additive} at the
  \(\Module \bar k\)-flavoured sheaf-cohomology level the project
  consumes; the additivity lemma is the project-side bridge.
\end{proof}

\section{The Riemann--Roch formula in genus zero}

\begin{theorem}


\leanok
[Riemann--Roch formula in genus zero]
  \label{thm:riemannRoch_genus_zero}
  \lean{AlgebraicGeometry.Scheme.WeilDivisor.l_eq_degree_plus_one_of_genus_zero}
  \uses{thm:euler_char_eq_deg_plus_one_minus_genus, def:eulerChar_curve,
        def:l_invariant}

  % SOURCE: [Hartshorne], IV.1, Example~1.3.5 (genus-\(0\) specialisation
  % of Riemann--Roch; the curve-is-rational consequence) (read from
  % references/hartshorne-algebraic-geometry.pdf, PDF page 314).
  % SOURCE QUOTE: "Example 1.3.5. Recall that a curve is \emph{rational}
  % if it is birational to \(\mathbf P^1\) (I, Ex. 6.1). Since curves in
  % this chapter are complete and nonsingular by definition, a curve
  % \(X\) is rational if and only if \(X \cong \mathbf P^1\) (I, 6.12). Now
  % using (1.3) we can show that \(X\) is rational if and only if \(g = 0\).
  % We already know that \(p_a(\mathbf P^1) = 0\) (I, Ex. 7.2), so
  % conversely suppose given a curve \(X\) of genus 0. Let \(P, Q\) be two
  % distinct points of \(X\) and apply Riemann--Roch to the divisor
  % \(D = P - Q\). Since \(\deg K - D = -2\), using (1.3.3) above, we have
  % \(l(K - D) = 0\), and so we find \(l(D) = 1\). But \(D\) is a divisor of
  % degree 0, so by (1.2) we have \(D \sim 0\), in other words $P \sim
  % Q\(. But this implies \)X$ is rational (II, 6.10.1)."
  \textit{Source: Hartshorne, IV.1, Example~1.3.5, p.~297.}
  Let \(C\) be a smooth proper geometrically irreducible curve over
  \(\bar k\) with \(g(C) = 0\), and let \(D \in \mathrm{Div}(C)\) with
  \(\deg(D) \geq 0\). Then
  \[
    \ell(D) \;=\; \deg(D) + 1.
  \]
\end{theorem}

\begin{proof}
  \leanok
  
  Specialising
  \cref{thm:euler_char_eq_deg_plus_one_minus_genus} to the genus-\(0\)
  hypothesis \(g = 0\) yields
  \[
    \chi(\mathcal O_C(D)) \;=\; \deg(D) + 1.
  \]
  Unfolding \(\chi\) via \cref{def:eulerChar_curve} and the definition
  \(\ell(D) = \dim_{\bar k} H^0(C, \mathcal O_C(D))\)
  (\cref{def:l_invariant}), this rewrites as
  \[
    \ell(D) - \dim_{\bar k} H^1(C, \mathcal O_C(D))
    \;=\; \deg(D) + 1.
  \]
  Under the standing genus-\(0\) hypothesis and the assumption
  \(\deg(D) \geq 0\), the higher-cohomology term vanishes:
  \[
    H^1(C, \mathcal O_C(D)) \;=\; 0
    \qquad \text{for } \deg(D) \geq 0.
  \]
  Substituting yields \(\ell(D) = \deg(D) + 1\), as required.

  \emph{The \(H^1\)-vanishing input.} On a smooth proper geometrically
  irreducible curve \(C / \bar k\) with \(g(C) = 0\), the vanishing
  \(H^1(C, \mathcal O_C(D)) = 0\) for \(\deg(D) \geq 0\) is the Hartshorne
  formulation Lemma~1.2 read covariantly: \(\ell(D) \neq 0\) forces
  \(\deg(D) \geq 0\); equivalently, for a curve where every
  degree-\(0\)-or-higher invertible sheaf is generated by global
  sections after a single point-twist, \(H^1\) of the twist must vanish
  by the long exact sequence of the skyscraper extension at any
  closed point. In the chain of sub-build chapters, this \(H^1\)
  vanishing is the natural content of \texttt{RR.3} (sibling chapter
  \texttt{RiemannRoch\_OcOfD.tex}, to be added), where the line
  bundle \(\mathcal O_C(D)\) and its cohomology are constructed; once
  \texttt{RR.3} is closed, the present theorem inherits the
  vanishing as an instance, and the proof closes verbatim as
  displayed. Pending \texttt{RR.3}, the Lean signature
  \texttt{l\_eq\_degree\_plus\_one\_of\_genus\_zero} threads the
  \(H^1\)-vanishing hypothesis explicitly as a named premise; the
  iteration that closes \texttt{RR.3} then drops the premise in a
  follow-up signature simplification.

  \emph{Sub-build note for \texttt{RR.4}.} The downstream consumer of
  this theorem is the genus-\(0\) classification
  \cref{prop:rigidity_genus0_curve_to_AV} of
  \texttt{AbelianVarietyRigidity.tex}: with \(\ell(D) = \deg(D) + 1\)
  and the degree-zero principal-divisor identity
  (\cref{thm:principal_deg_zero}), Hartshorne's Example~1.3.5 chain
  produces a degree-\(1\) morphism \(C \to \mathbb P^1_{\bar k}\) via
  \texttt{Proj.fromOfGlobalSections} applied to the two-dimensional
  \(H^0(C, \mathcal O_C([P]))\), which is then upgraded to an
  isomorphism in \texttt{RR.4}.
\end{proof}

\section{Out of scope}

The following items are deliberately \emph{not} part of \texttt{RR.2}
and are sequenced into the sibling sub-build chapters:

\begin{itemize}
  \item The general-\(g\) Riemann--Roch theorem
    \(\ell(D) - \ell(K - D) = \deg(D) + 1 - g\) (Hartshorne~IV.1.3 in
    its full form). This requires the canonical divisor \(K\), the
    canonical sheaf \(\Omega^1_{C / \bar k} = \omega_C\), and
    Serre duality $H^1(C, \mathcal L) \cong H^0(C, \omega_C \otimes
    \mathcal L^{-1})^*$. Mathlib has none of these for schemes, and
    the genus-\(0\) critical path does not require them; the general
    statement is queued for a follow-up chapter (post-RR-bridge).
  \item The construction of the invertible sheaf \(\mathcal O_C(D)\)
    from a Weil divisor \(D\), the linear-equivalence isomorphism
    \(\mathcal O_C(D) \cong \mathcal O_C(D')\) when \(D \sim D'\), and
    the computation of $H^0(C, \mathcal O_C([P])) \cong \bar k \cdot
    1 \oplus \bar k \cdot t$ (\emph{a non-constant rational function
    \(t\) exists}) live in \texttt{RR.3} (sibling chapter
    \texttt{RiemannRoch\_OcOfD.tex}, to be added). Inside the proof
    of \cref{thm:euler_char_eq_deg_plus_one_minus_genus} the dependency on
    \(\mathcal O_C(D)\) is reduced to its \(H^0\) functor and the short
    exact sequence of the closed point.
  \item The vanishing \(H^1(C, \mathcal O_C(D)) = 0\) for
    \(\deg(D) \geq 0\) on a genus-\(0\) curve is consumed as a named
    premise by \cref{thm:riemannRoch_genus_zero}; its closure is
    part of \texttt{RR.3} (where the cohomology of
    \(\mathcal O_{\mathbb P^1}(d)\) is computed) and is recovered for
    a general smooth proper genus-\(0\) curve via the isomorphism
    \(C \cong \mathbb P^1\) of \texttt{RR.4} or directly via the
    base-change argument inside \texttt{RR.3}.
  \item The classification ``a smooth proper geometrically irreducible
    curve over \(\bar k\) with \(g = 0\) is isomorphic to
    \(\mathbb P^1\)'' (Hartshorne Example~IV.1.3.5 proper) is the
    headline of \texttt{RR.4} (sibling chapter
    \texttt{RiemannRoch\_RationalIsoP1.tex}, to be added). The
    present chapter supplies one of its two inputs (the $\ell(D) =
    \deg(D) + 1$ identity); the other input is the degree-$0$
    principal-divisor identity \cref{thm:principal_deg_zero} from
    \texttt{RR.1}.
  \item Higher-cohomology bookkeeping (\(\chi\) as an alternating sum
    over all \(i \geq 0\), the Grothendieck vanishing
    \(H^i(C, \mathcal F) = 0\) for \(i > \dim C = 1\)) is not separately
    formalised here; the chapter consumes only the two-term form
    \(\chi(\mathcal F) = \dim H^0 - \dim H^1\) which is the standing
    curve specialisation.
\end{itemize}
